\chapter{Uncertainty, Hilbert Space, and the Oscillator}
\label{chap:uncertainty-hilbert-oscillator}

Part IV is the book's largest, and it opens through a door the last part left ajar. Brownian motion showed
that the diffusion of a finite random walk, continued a quarter-turn into imaginary time, is the
Schr\"odinger evolution --- so the quantum is diffusion run in imaginary time, and the count the instrument
has carried since the first mark is carried now as an amplitude, a complex number whose square is a
probability. Everything in this part is that one change rung over and over: the count becomes an amplitude.
This first chapter lays the three foundations everything quantum rests on --- the uncertainty that fixes the
smallest cell the instrument can carve, the Hilbert space the amplitudes live in, and the harmonic oscillator
whose quantized ladder is the simplest quantum spectrum --- and two of the three come in two parts, because
each carries two distinct readings the instrument keeps separate.

\section{The Heisenberg Effect, I: the floor}
\label{se:heisenberg-floor}

Heisenberg's inequality,
\begin{equation}
  \Delta x\,\Delta p \ge \frac{\hbar}{2},
\end{equation}
has two readings, and the instrument keeps them apart because they stand on two different honest floors. The
first is a floor in the literal sense: a smallest joint cell. No measurement, however fine the instrument,
can pin a particle's position and momentum together more tightly than $\hbar/2$ --- there is a minimum joint
refinement below which the record simply cannot go. To the instrument this is not a statement about
disturbance or clumsiness but about the size of the smallest cell it can carve in the joint space of position
and momentum: $\hbar/2$ is the area of that cell, the finest joint distinction the record admits. Sharpen the
position to a point and the momentum cell stretches to fill the whole axis; the product of the two spreads
holds at $\hbar/2$, the floor.

Why does the floor exist at all? Beneath the inequality lies a deeper fact, and it is the root of the whole
chapter: position and momentum do not commute. Apply them to a record in one order and then the other, and
the two results differ by a fixed amount,
\begin{equation}
  [\,x, p\,] = x p - p x = i\hbar .
\end{equation}
This is the canonical commutator, and it says that $x$ and $p$ are not two independent dials the instrument
can set apart but conjugate quantities, each the other's Fourier transform: a sharp position is a wave packet
assembled from all momenta, a sharp momentum a wave spread across all positions, and the two can never both be
sharp because they are two views of one object. The uncertainty inequality is the shadow this non-commutation
casts. In its general form --- the Robertson inequality --- the spreads of any two quantities are bounded
below by the size of their commutator,
\begin{equation}
  \Delta A\,\Delta B \;\ge\; \tfrac12\bigl|\langle [\,A, B\,]\rangle\bigr|,
\end{equation}
and for position and momentum the right-hand side is $\tfrac12|i\hbar| = \hbar/2$. The floor is not an extra
postulate; it is the direct consequence of $[x,p] = i\hbar$. The smallest joint cell exists because $x$ and
$p$ refuse to commute, and its area $\hbar/2$ is the magnitude of their irreducible interference. An unbounded
joint sharpening would demand that the two commute after all --- that an infinite summation of finite
refinements pin a point in a continuous joint plane --- and the commutator forbids it.

The honest floor is a finite count obligation: $\hbar/2$ is the area of the smallest joint cell, a count the
record cannot subdivide. The reading is that uncertainty, read as a floor, is the resolution of the conjugate
plane --- the finest joint distinction the instrument can draw, set by the finiteness of its records and not
by any clumsiness of measurement.

\section{The Heisenberg Effect, II: the trade-off}
\label{se:heisenberg-tradeoff}

The second reading of the same inequality is a trade-off, and it stands on a different floor. Read this way,
$\hbar/2$ is not the size of a fixed cell but the total of a fixed budget --- a measurement-capacity the
instrument has to spend, and must divide between the two conjugates. Spend it on the position, sharpening
$\Delta x$, and there is less left for the momentum, so $\Delta p$ must grow; spend it on the momentum and the
position blurs. The product $\Delta x\,\Delta p$ is the budget, and it cannot be beaten: there is no way to
buy a sharper position without paying in a blurrier momentum, no way to get both for less than $\hbar/2$.

The trade-off is visible in the simplest experiment. Send a beam through a single slit and narrow it: as the
slit closes, the position of each particle passing through is pinned more tightly --- $\Delta x$ shrinks ---
and the beam, instead of staying a thin line, fans out behind the slit into a diffraction pattern, wider the
narrower the slit. That spreading is $\Delta p$ growing to pay for the sharpened $\Delta x$. Close the slit
toward a point and the beam fans across the whole screen; open it wide and the beam stays narrow but its
position is vague. No setting gives both a sharp position and a sharp direction, because the slit is spending
the one budget the inequality names.

And here the floor becomes a wall. Driving the product below $\hbar/2$ is not merely hard or expensive --- it
is inadmissible. There is no record at all with $\Delta x\,\Delta p < \hbar/2$; the instrument can carve the
cell tall-and-thin or short-and-wide, but it cannot carve it smaller than $\hbar/2$, the same way it cannot
draw a slip below the friction threshold or a transmission through crossed polarizers. The honest floor here
is a no-go, the trade-off's wall, distinct from the finite-count floor of the previous section: the floor
says the cell has a smallest size; the wall says no record can have a cell smaller still.

The reading is that uncertainty is not ignorance but a capacity. The instrument has a finite budget for
distinguishing conjugates, and it spends that budget where it chooses; what it cannot do is spend less than
the budget. The trade-off is the budget, and the floor is its edge --- two faces of one inequality, a count
and a wall, which is why the chapter gives Heisenberg two sections and not one.

\section{The Hilbert Effect: the space the amplitudes live in}
\label{se:hilbert}

If the count is now an amplitude, the amplitudes need somewhere to live, and the instrument's own refinements
supply it. The admissible refinements do not stay a loose collection of states; completed under prediction
and consistency, they form an inner-product space --- a space with a notion of length and angle. Distances
obey a Pythagorean norm,
\begin{equation}
  \lVert p\rVert^2 = a^2 + b^2,
\end{equation}
so that the size of an amplitude is got from its parts exactly as a hypotenuse is got from its legs;
amplitudes combine through that inner product, which measures how much one shares with another; and prediction
is carried by an operator $\Psi$ that has an inverse $\Psi^{-1}$, the two mutual inverses, so the evolution
can be run forward and backward alike --- the quantum is reversible, unlike the diffusion it is the
imaginary-time cousin of.

The inner product does the real work, and it pays to write it out. Taking two states $\psi$ and $\phi$, it
returns a number,
\begin{equation}
  \langle \psi \mid \phi \rangle,
\end{equation}
the overlap of the one with the other --- the amplitude to find a system prepared as $\phi$ when the
instrument tests it for $\psi$. Its square is the probability, so the geometry of the space is the physics of
the chances: two states that are orthogonal, $\langle \psi \mid \phi \rangle = 0$, exclude one another, while
a state's overlap with itself, $\langle \psi \mid \psi \rangle = \lVert \psi \rVert^2 = 1$, is the certainty
that it is what it is. States add --- superpose --- to make new states, and a measurement's chances are read
off the angles between them.

The quantities the instrument can measure are operators on this space, and not arbitrary ones: an observable
is an operator whose eigenvalues are real, because a measurement returns a real number. An observable $A$ has
eigenstates $\lvert a\rangle$ and eigenvalues $a$,
\begin{equation}
  A\,\lvert a\rangle = a\,\lvert a\rangle,
\end{equation}
and the eigenvalues are exactly the values the measurement can return, the eigenstates the records in which
the observable has a definite value. Any state expands over an observable's eigenstates, $\lvert \psi\rangle
= \sum_a c_a \lvert a\rangle$, and the squared coefficient $|c_a|^2$ is the probability the measurement
returns $a$ --- the spectral picture, where the possible outcomes are the spectrum of the operator and their
probabilities are the squared overlaps. Position and momentum are two such observables, and the reason they
cannot both be sharp now reads off the structure directly: they share no eigenstates, because their
commutator $[x,p] = i\hbar$ is not zero, and a state cannot be an eigenstate of both.

The last ingredient is completeness, and it must be said carefully, in plain words. The space is complete:
every sequence of admissible refinements that ought to converge does converge, to a limit that is itself an
admissible refinement. The record is closed under the limits it can take --- it has no holes, no missing
points where a convergent sequence runs off the edge of the space. A complete inner-product space is, by
definition, a Hilbert space, and so the amplitudes' home is a Hilbert space not by decree but by completion:
take the instrument's finite inner-product structure, close it under its own limits, and a Hilbert space is
what you have.

The honest floor is a smooth-shadow analogy: the discrete inner-product structure of the refinements is the
model, the continuum Hilbert space its smooth shadow, and the completion --- the filling-in of the limits ---
is the bridge, named as one. The reading is that quantum amplitudes live in a Hilbert space because that is
what the instrument's own record becomes when it is closed under its limits: reversibility is the prediction
operator having an inverse, and completeness is the record having no holes.

\section{The Harmonic Oscillator, I: the dynamics}
\label{se:oscillator-dynamics}

The harmonic oscillator is the simplest system the instrument can build on a conjugate pair, and it is worth
meeting in two stages: first its dynamics, the reversible swing, and then its spectrum, the quantized ladder.
The dynamics are a reciprocity --- a steady trading of position for momentum and back, each feeding the
other,
\begin{equation}
  \dot{x} = \beta p, \qquad \dot{p} = -\alpha x,
\end{equation}
the infinitesimal form of a rotation in the conjugate plane, $(x, p) \mapsto (x + \beta p,\, p - \alpha x)$:
a little step turns a bit of position into momentum and a bit of momentum into position, and iterated, the
point $(x, p)$ circles the origin forever. In the continuum this circling is simple harmonic motion,
\begin{equation}
  \delta^2 U + \omega^2 U = 0 \;\longrightarrow\; \frac{d^2 U}{dt^2} + \omega^2 U = 0,
\end{equation}
the equation of a swing, with a conserved energy
\begin{equation}
  E = \tfrac12\!\left[ (\dot U)^2 + \omega^2 U^2 \right]
\end{equation}
that the rotation holds fixed --- the radius of the circle in the conjugate plane never changes.

The honest floor is a smooth-shadow analogy: the discrete reciprocity --- the step-by-step rotation of the
conjugate pair --- is the model, and the smooth harmonic equation $d^2 U/dt^2 + \omega^2 U = 0$ is its shadow,
the continuum the discrete circling approaches. The reading is that the oscillator's swing is reversible
rotation in the conjugate plane: not a force pulling a mass back to center, but a steady reciprocity trading
the two conjugates, around and around, at a fixed radius.

\section{The Harmonic Oscillator, II: the ladder}
\label{se:oscillator-ladder}

Quantize that swing and its energy stops being continuous. Where a classical oscillator can circle at any
radius, holding any energy, the quantum oscillator's energy comes in a ladder of evenly spaced rungs,
\begin{equation}
  E_n = \hbar\omega\left(n + \tfrac12\right),
\end{equation}
each rung one quantum $\hbar\omega$ above the last, $n$ counting the rungs from the bottom. The spacing is the
finite count made visible: the oscillator's energy is a whole number of quanta, $n$ of them, and nothing
between --- the count, which became an amplitude at the chapter's start, is here a rung number.

The ladder is not a metaphor; the instrument climbs it with two operators built from the conjugate pair.
Combine position and momentum into a complex pair,
\begin{equation}
  a = \tfrac{1}{\sqrt2}\,(x + i p), \qquad a^\dagger = \tfrac{1}{\sqrt2}\,(x - i p),
\end{equation}
the lowering and raising operators, each the other's conjugate, with a commutator that inherits the
canonical one,
\begin{equation}
  [\,a, a^\dagger\,] = 1 .
\end{equation}
From them comes the number operator $N = a^\dagger a$, whose eigenvalues are exactly the rung numbers,
\begin{equation}
  N\,\lvert n\rangle = n\,\lvert n\rangle,
\end{equation}
with $n$ a non-negative integer counting the quanta in the state. And the names are earned. The raising
operator carries the oscillator up one rung, the lowering operator down,
\begin{equation}
  a^\dagger\,\lvert n\rangle = \sqrt{n+1}\,\lvert n+1\rangle, \qquad
  a\,\lvert n\rangle = \sqrt{n}\,\lvert n-1\rangle .
\end{equation}
Read through these operators, the energy $E_n = \hbar\omega(n + \tfrac12)$ is the instrument climbing and
descending a ladder one rung at a time, and each rung is one count: $a^\dagger$ is the act of adding a single
quantum to the tally, $a$ the act of removing one. The amplitude becomes a count again here, in the most
literal way --- the state $\lvert n\rangle$ \emph{is} $n$ quanta, and the operators that build and unbuild it
step the count by exactly one.

But look at the bottom rung. It is not zero. Even at $n = 0$, the lowest state, the oscillator holds energy
$\tfrac12\hbar\omega$ --- the zero-point --- and it cannot give it up. This is the cross-link the chapter has
been building toward, the place where the oscillator's last section meets Heisenberg's first. The oscillator
cannot sit still, cannot rest at the bottom of its well with both position and momentum exactly zero, because
a state of zero spread in both $x$ and $p$ would be a joint cell of zero area, and the uncertainty floor
forbids it: $\Delta x\,\Delta p \ge \hbar/2$. Algebraically, the floor is the lowering operator running out of
rungs: applied to the ground state it returns nothing,
\begin{equation}
  a\,\lvert 0\rangle = 0,
\end{equation}
because there is no rung below the lowest, and $a\lvert 0\rangle = 0$ is the exact algebraic statement of ``the
oscillator cannot sit still.'' The smallest the oscillator can do is occupy the smallest cell the floor
allows, and the energy of that smallest cell is exactly the zero-point $\tfrac12\hbar\omega$. The
zero-point is the uncertainty floor, given a ladder; the floor that the first section called the smallest
joint cell is, here, the lowest rung the oscillator can stand on.

The zero-point is not a bookkeeping fiction; it is measured. Helium, cooled toward absolute zero, refuses to
freeze under its own vapor pressure --- its atoms jitter with a zero-point motion too restless to lock into a
solid, and only squeezing it hard enough to overcome that motion makes it crystallize. The same zero-point
restlessness, in the modes of the electromagnetic field, is what the Casimir effect of the heat part measured
--- two uncharged plates drawn together by the vacuum's own jitter. The smallest cell the uncertainty floor
allows is a real, restless energy, and the bottom of the oscillator's ladder is where the instrument first
holds it as a number.

The honest floor is a finite count obligation: the spectrum is a whole number of quanta above the zero-point,
the count being the rung number $n$. The reading is that the oscillator is the simplest place the count
becomes an amplitude and the amplitude a ladder --- and the bottom of that ladder is the uncertainty floor
itself, the smallest cell a finite record can carve, lifted into the lowest energy a quantum oscillator can
hold. The two un-folded effects of this chapter, the floor and the ladder, are one fact read at its two ends.

\bigskip

With the floor, the space, and the ladder, Part IV has its foundations, and the count has become an amplitude
in three ways. Uncertainty fixes the smallest cell the instrument can carve in the conjugate plane --- a count
for a floor, a wall against carving smaller. The Hilbert space gives the amplitudes a complete home, the
instrument's own record closed under its limits. And the oscillator turns the floor into a spectrum, its
lowest rung the smallest cell lifted into an energy. The part takes up next what these foundations make
strange: measurement, decoherence, and the nonlocal correlations the instrument reads when amplitudes, not
counts, are what cross a boundary.
